\documentclass[]{article}

%Packages________________
\usepackage[normalem]{ulem}
\usepackage[dvipsnames]{xcolor}
\usepackage{cancel}
\usepackage{parskip}
\usepackage{caption, subcaption}
\usepackage[Spanish,es-tabla,es-noshorthands]{babel}
\usepackage{newtxtext}
\usepackage[varvw]{newtxmath}
\usepackage{booktabs, colortbl, bigstrut, multirow, multicol}
\usepackage{url}
\usepackage{pgfplots}
\usepackage{graphicx}
\usepackage{amsmath}
\usepackage{tabularx}
\usepackage{floatrow}
\usepackage{float}
\usepackage{cancel}
\usepackage[spanish]{cleveref}
\usepackage{tikz}
\usetikzlibrary{arrows.meta,shapes.geometric,positioning,calc}
\crefname{table}{tabla}{tablas}
\usepackage{wrapfig}
\usepackage{adjustbox}
\usepackage{geometry}
\usepackage{lipsum}
\usepackage{fancybox}
\usepackage[T1]{fontenc}
\geometry{
	top=2.4cm,
	bottom=2.4cm,
	left=2.3cm,
	right=2.3cm
}

\usepackage[
    backend=biber,
	citestyle=verbose-ibid,
    bibstyle=authoryear-icomp,
    natbib=true,
    url=true, 
    doi=true,
    eprint=false
]{biblatex}

\addbibresource{ref.bib}

%Commands________________
\newcommand\hcancel[2][black]{\setbox0=\hbox{$#2$}%strikeout
\rlap{\raisebox{.45\ht0}{\textcolor{#1}{\rule{\wd0}{1pt}}}}#2} 

\renewcommand\theequation{\Alph{equation}}

\def\mydate{\leavevmode\hbox{\twodigits\day/\twodigits\month/\the\year}}
\def\twodigits#1{\ifnum#1<10 0\fi\the#1}

\newcommand\runinsec{\@startsection {section}{1}{\z@}%
                                   {3.5ex \@plus 1ex \@minus .2ex}%
                                   {-1em}%
                                   {\normalfont\Large\bfseries}}
\newcommand\aftersec[1]{\begingroup\normalfont\footnotesize #1\endgroup%
                        \par\nobreak\vspace*{2.3ex}\noindent\ignorespaces}

%\renewcommand\thesubsection{\alph{subsection}}

% \input{colorbox.txt}

%Titlepage_________________
\title{\vspace{-1.8cm} Title}
\author{Álvaro Jerónimo Sánchez}
\date{}

%Document_____________
\begin{document}
\maketitle
\setcounter{section}{0}
\renewcommand{\tablename}{Tabla}


\section{Ejercicio 1}

\subsection{Matriz}

Los mecanismos de \textit{reacción-difusión}, propuestos por Alan Turing, explican cómo los patrones presentes en diversos animales son resultado de reacciones entre dos o más sustancias químicas (activador-inhibidor) con pequeñas perturbaciones que acaban causando estructuras periódicas. En este ejercicio se pretende resolver las ecuaciones de dicho modelo creando una red $L \times L$, introduciendo perturbaciones y calculando el \textit{laplaciano} en cada punto discreto; todo a partir de unos parámetros iniciales dados.

(...)

Se ha generado el ruido gaussiano aplicando la transformada de Box-Muller, cuya fórmula es: 

\subsection{Amplitud}

Para hallar la amplitud, se calcula el promedio de $(u-u0)^2$ para cada timestep, siendo u el valor de cada punto de la red en ese timestep, y u0 el valor fijo inicial sin las perturbaciones (a+b)

%_________________________________________

%\nocite{*}
%\printbibliography

\appendix
\textbf{APPENDIX}

\end{document}